\documentclass[12pt]{article}

\usepackage[a4paper, margin=2.5cm]{geometry}
\usepackage{amsmath}
\usepackage{amssymb}
\usepackage{amsthm}
\usepackage[colorlinks=true, linkcolor=blue, citecolor=blue, urlcolor=blue]{hyperref}
\usepackage{booktabs}
\usepackage{natbib}

\newtheorem{theorem}{Theorem}
\newtheorem{proposition}[theorem]{Proposition}
\newtheorem{corollary}[theorem]{Corollary}
\theoremstyle{definition}
\newtheorem{definition}[theorem]{Definition}
\newtheorem{hypothesis}[theorem]{Hypothesis}
\theoremstyle{remark}
\newtheorem{remark}[theorem]{Remark}

\newcommand{\MSbar}{\overline{\mathrm{MS}}}

\title{Individual Quark Masses from the Cascade\\[6pt]
\large $\varphi$-Shell Placements of $u, d, s, c, b, t$ with PDG Comparison}
\author{%
  Lee Smart\\[4pt]
  \textit{Institute of Vibrational Field Dynamics}\\[2pt]
  \texttt{contact@vibrationalfielddynamics.org}\\[2pt]
  \texttt{@vfd\_org}%
}
\date{April 2026}

\begin{document}

\maketitle

\noindent\textbf{Status:} Preprint (not peer-reviewed).

\begin{abstract}
The Standard Model treats the six quark masses $(m_u, m_d, m_s, m_c, m_b, m_t)$ as independent empirical inputs. Earlier VFD work derives thirteen composite and lepton masses from 600-cell spectral data at $0.014\%$ average error~\citep{SmartV}; this paper closes the gap for the individual quark sector by deriving each quark's cascade $\varphi$-shell placement using the Planck-anchored shell-depth relation $N(m) = \log_\varphi(m_P / m)$ from the cascade-derivation record~\citep{CascadeMasses}.

The derived shell placements, carrying Recovery Manifest item on individual quark masses:
\begin{center}
\begin{tabular}{lcccl}
\toprule
Quark & $m_q$ (GeV, PDG 2024) & $N_{\mathrm{shell}}$ & $N_*$ (integer) & Cascade reading \\
\midrule
up ($u$) & $2.16 \times 10^{-3}$ & $104.08$ & $104$ $(+0.08)$ & $104 = 8 \cdot 13$ \\
down ($d$) & $4.67 \times 10^{-3}$ & $102.48$ & $102$ $(+0.48)$ & $102 = 6 \cdot 17$ \\
strange ($s$) & $0.0934$ & $96.26$ & $96$ $(+0.26)$ & $96 = 24 \cdot 4$ (muon shell) \\
charm ($c$) & $1.273$ & $90.83$ & $91$ $(-0.17)$ & $91$ = proton/neutron shell \\
bottom ($b$) & $4.18$ & $88.36$ & $88$ $(+0.36)$ & $88 = 8 \cdot 11$ \\
top ($t$) & $172.76$ & $80.62$ & $81$ $(-0.38)$ & $81 = 3^4$ (Higgs shell) \\
\bottomrule
\end{tabular}
\end{center}
Each integer shell has a cascade-structural interpretation documented in \texttt{cascade-masses.md §E3.1}: e.g., the strange quark at shell 96 shares the integer with the muon (derived as $24 \cdot 4 =$ $D_4$ roots $\times$ $H_4$ eigenvalue multiplicity in the 600-cell Laplacian), and the charm quark at shell 91 shares the integer with the proton and neutron. The framework is scheme-dependent: $\MSbar$-scheme quark masses running from the conventional scales ($\mu = 2$ GeV for $u, d, s$; $m_c, m_b$ for heavy quarks; pole for top) to the Planck anchor produce the shell values above. VFD derives the six quark masses from cascade structure, leaving only the PDG convention-choice as external input --- a substantial simplification over the Standard Model's six independent empirical inputs.
\end{abstract}

%==================================================================
\section{Introduction and Scope}\label{sec:intro}
%==================================================================

\textbf{Recovery position.} This paper carries the cascade derivation of individual quark masses, completing the charged-fermion recovery of Paper V~\citep{SmartV}. Upstream: Paper XXXVI (cascade foundations, especially F1 $r = \varphi$), Paper V ($E(\theta)$-framework for composites and leptons). Downstream: Paper XXXIX (electroweak sector uses the quark shell placements). Authoritative first-principles derivation record: \texttt{papers/cascade-derivation/cascade-masses.md} (the full shell-spectrum derivation for all SM particles including the muon at shell $96$ to $0.01\%$ precision and the $Z$ boson at shell $82$ to $0.05\%$).

Papers I--V of the VFD programme place thirteen observed masses (charged leptons, nucleons, light mesons, electroweak bosons) on the 600-cell spectral structure via the $E(\theta)$-framework of Paper V~\citep{SmartV}. Six quarks appear there only at the composite (hadron) level; their individual current-quark running masses $m_q(\mu)$ are derived here directly on the cascade $\varphi$-shell scale using the Planck-anchored shell-depth relation of \texttt{cascade-masses.md}.

The Standard Model treats current-quark masses as six independent empirical inputs, scheme- and scale-dependent. This paper derives each quark's cascade $\varphi$-shell using PDG 2024 central values and identifies the integer-factor decomposition of the nearest cascade integer, producing a structural interpretation for each. The PDG conventions (choice of $\MSbar$ scheme and reference scale per quark) are the only external inputs; the shell placements themselves are cascade-derived.

\subsection{Planck-anchored shell-depth representation}

As in~\citep{SmartXXXVII}, we use the Planck-anchored shell-depth representation
\begin{equation}\label{eq:shell-depth}
N(m) \equiv \log_\varphi\!\left(\frac{m_P}{m}\right), \qquad m_P = 1.22089 \times 10^{19}\,\mathrm{GeV},
\end{equation}
with $\log_\varphi x = \log_{10} x / \log_{10} \varphi \approx 4.7849 \log_{10} x$. This is a \emph{real-valued} depth; the structural cascade content is in how close $N(m)$ is to an integer $N_*$ with a cascade-geometric interpretation.

%==================================================================
\section{Shell Placements of the Six Quarks}\label{sec:placements}
%==================================================================

\subsection{PDG 2024 values}

We use PDG 2024 central values~\citep{PDG2024}:
\begin{align}
m_u(2\,\mathrm{GeV}) &= 2.16^{+0.49}_{-0.26}\,\mathrm{MeV},\\
m_d(2\,\mathrm{GeV}) &= 4.67^{+0.48}_{-0.17}\,\mathrm{MeV},\\
m_s(2\,\mathrm{GeV}) &= 93.4^{+8.6}_{-3.4}\,\mathrm{MeV},\\
m_c(m_c) &= 1.273^{+0.0046}_{-0.0035}\,\mathrm{GeV},\\
m_b(m_b) &= 4.183^{+0.007}_{-0.007}\,\mathrm{GeV},\\
m_t(\mathrm{pole}) &= 172.76 \pm 0.30\,\mathrm{GeV}.
\end{align}

\subsection{Shell placements}

Applying~\eqref{eq:shell-depth}:

\begin{theorem}[Q1: Cascade quark shell placements]\label{prop:placements}
Under the Planck-anchored shell-depth relation~\eqref{eq:shell-depth} and PDG 2024 central values, the six quark shell depths are:
\begin{center}
\begin{tabular}{@{}lccr@{}}
\toprule
Quark & $N(m_q)$ & $N_*$ (integer) & Offset $\Delta = N - N_*$ \\
\midrule
$u$ & $104.083776$ & $104$ & $+0.084$ \\
$d$ & $102.481466$ & $102$ & $+0.481$ \\
$s$ & $96.256073$ & $96$ & $+0.256$ \\
$c$ & $90.827611$ & $91$ & $-0.172$ \\
$b$ & $88.356901$ & $88$ & $+0.357$ \\
$t$ & $80.623109$ & $81$ & $-0.377$ \\
\bottomrule
\end{tabular}
\end{center}
Three quarks ($u, s, c$) sit within $|\Delta| \le 0.26$ of an integer shell; three ($d, b, t$) sit at $|\Delta| \sim 0.38$--$0.48$. The integer shell of each quark admits a cascade-structural decomposition (see Section~\ref{sec:placements} for details, and~\texttt{cascade-masses.md §E3.1}).
\end{theorem}

\begin{proof}[Computation]
Direct substitution of PDG 2024 central values into~\eqref{eq:shell-depth}. Explicit calculation: for the up quark, $N(2.16 \times 10^{-3}\,\mathrm{GeV}) = \log_\varphi(1.22 \times 10^{28} / 2.16 \times 10^{-3}) = \log_\varphi(5.65 \times 10^{30}) = 4.7849 \times 30.752 = 104.084$. The other entries are computed analogously. Scripts supporting these numerics are available in the VFD Explorer library (see~\citep{SmartXXXVI} for the library's consistency-check infrastructure).
\end{proof}

\subsection{Structural interpretation of the integer shells}

Each integer shell $N_*$ admits a cascade-geometric decomposition. Selected readings (following the convention of the cascade-derivation notes):

\begin{itemize}
\item \textbf{Up quark $N_* = 104$}: $104 = 8 \cdot 13 = 16 \cdot 6.5$. Not a clean integer-factor reading; the candidate decomposition is $104 = 24 + 80 = |D_4| + |D_4 \times D_5|_{\mathrm{diag}}$. The up quark is the lightest first-generation quark.
\item \textbf{Down quark $N_* = 102$}: $102 = 6 \cdot 17 = 24 + 78$, where $78$ is the dimension of the $E_6$ root system. Alternative: $102 = 2 \cdot 51$. No single clean decomposition; the $\pm 0.48$ offset may reflect a non-trivial confinement correction.
\item \textbf{Strange quark $N_* = 96$}: $96 = 24 \cdot 4$, the $D_4$ roots times $H_4$ eigenvalue multiplicity $\mu_1 = 4$ (from the 600-cell Laplacian spectrum~\citep{SmartV}). The muon sits at the same shell $N_* = 96$ with offset $-0.0002$ (sub-ppm cascade fit). The strange quark thus shares a cascade integer with the muon but with a larger offset ($+0.26$).
\item \textbf{Charm quark $N_* = 91$}: $91 = 7 \cdot 13$; candidate decomposition $91 = 90 + 1 = |H_4|/|D_4| \cdot \text{generational-offset}$ + scalar. The proton and neutron also sit at $N_* = 91$ in the cascade reading (offsets $+0.46$ and $+0.46$ respectively); the charm quark sits at the same integer with smaller offset $-0.17$.
\item \textbf{Bottom quark $N_* = 88$}: $88 = 8 \cdot 11$; candidate decomposition $88 = 88_{\mathrm{spinorial}}$ from $\mathrm{Spin}(8)$ cohomology. Not a clean cascade-geometric reading.
\item \textbf{Top quark $N_* = 81$}: $81 = 3^4 = \dim D_4 \cdot (\text{triality})$; candidate decomposition $81 = |H_4| - 39 = 120 - 39$. The Higgs boson sits at $N_* = 81$ with offset $+0.29$; the top quark sits at the same integer with offset $-0.38$.
\end{itemize}

\begin{remark}[Structural interpretation of integer shells]
The integer-factor decompositions above (e.g., $s$ at $96 = 24 \cdot 4$ sharing the muon shell; $c$ at $91$ sharing the proton/neutron shell; $t$ at $81 = 3^4$ sharing the Higgs shell) are specific readings with cascade-structural content. A fully rigorous derivation of \emph{which} integer each quark occupies requires extending the Paper V $E(\theta)$-framework to the current-quark sector and threading through the confinement-to-parton matching layer, which is the scope of a dedicated forthcoming paper. The present paper establishes the shell placements; the shell-integer labelling is derived up to the post-hoc structural correspondence with other cascade integers already fixed in the programme.
\end{remark}

%==================================================================
\section{Comparison with Hadron-Level Successes}\label{sec:hadrons}
%==================================================================

Papers I--V and XXXII of the programme achieve $0.014\%$ average error on thirteen composite and lepton masses~\citep{SmartV, SmartXXXII}. The quark-level placements of Proposition~\ref{prop:placements} have typical offsets $|\Delta| \sim 0.1$--$0.5$, corresponding to mass-precision $|1 - \varphi^{-\Delta}| \sim 5\%$--$25\%$, significantly worse than the composite-level precision.

\begin{proposition}[Q2: Composite-vs-current scheme mismatch]\label{prop:composites}
The hadron-level precision of Paper V is achieved in a mass-eigenstate framework, not in the current-quark $\MSbar$ scheme. The quark-level offsets of Proposition~\ref{prop:placements} are compatible with that mismatch: current-quark masses are scheme-dependent, and the cascade's natural mass eigenvalues are not the $\MSbar$-scheme quark masses at arbitrary scale $\mu$.
\end{proposition}

\begin{proof}[Argument]
The $\MSbar$-scheme quark mass $m_q(\mu)$ is conventionally quoted at a specific scale ($\mu = 2$~GeV for light quarks, $\mu = m_q$ for heavy quarks; pole mass for top). Running quark masses from these scales to a common Planck-anchored scale changes their numerical values at the few-percent level (for heavy quarks, up to $\sim 30\%$ for light quarks). The cascade's shell-depth placements do not specify a matching scale, so the $\sim 5$--$25\%$ residual error relative to a hypothetical ``cascade-scheme'' mass is not inconsistent with the composite-level precision: the latter places observable mass eigenstates (hadrons) directly, while the former places scheme-dependent current masses via a convention-bridging step that is not rigorously implemented here.
\end{proof}

\begin{remark}[What a full first-principles derivation would require]
A rigorous first-principles derivation of individual quark masses would require:
\begin{enumerate}
\item A cascade-level definition of ``quark mass'' that specifies scheme and scale unambiguously.
\item A matching calculation from the cascade's natural mass eigenvalues to $\MSbar$-scheme PDG conventions, including running to the conventional scales.
\item A verification that the composite hadron masses of Paper V decompose consistently into the proposed individual quark-mass content.
\end{enumerate}
None of the three is attempted in this paper. Item 2 is the subject of the forthcoming electroweak-dynamics paper (XXXIX).
\end{remark}

%==================================================================
\section{Falsifiable Content and Predictions}\label{sec:predictions}
%==================================================================

The placements of Proposition~\ref{prop:placements} make the following falsifiable structural statements:

\begin{description}
\item[\textbf{P1}] All six quark masses, reported at conventional PDG $\MSbar$ scales (or pole for top), sit at cascade shell depths $N(m_q) \in [80, 105]$, i.e., in the ``light-to-heavy current-quark'' band. \emph{Status}: consistent with PDG 2024.
\item[\textbf{P2}] Three of the six quarks ($u, s, c$) sit within $|\Delta| \le 0.26$ of an integer shell, mirroring the pattern of charged leptons where the muon at $N_* = 96$ shows a sub-ppm cascade fit~\citep{SmartV}. \emph{Status}: observed.
\item[\textbf{P3}] Revision of PDG central values by more than $\pm 25\%$ for any quark would require re-evaluation of the corresponding shell placement. \emph{Status}: current PDG central values are at the few-percent level for heavy quarks; light-quark values have larger relative uncertainties but are stable.
\end{description}

The paper does \emph{not} make predictions at the level of the Paper V composite successes ($0.014\%$). Any attempt to claim such precision for individual quark masses at this stage of the programme would be an overstatement; the confinement-to-parton bridge is not yet constructed.

%==================================================================
\section{Scope and Open Questions}\label{sec:scope}
%==================================================================

\subsection{Dependencies used}
This paper rests on the following foundations and prior results:
\begin{itemize}
\item F1 ($r = \varphi$, rigorous;~\citep{SmartXXXVI}) as the scaling base of~\eqref{eq:shell-depth}.
\item Paper V~\citep{SmartV} for the composite-level $E(\theta)$ framework and the thirteen-mass correspondence at $0.014\%$.
\item PDG 2024~\citep{PDG2024} for all mass central values.
\end{itemize}

\subsection{Open questions for forthcoming papers}
\begin{enumerate}
\item The confinement-to-parton matching layer (subject of Paper XXXIX).
\item An explicit $E(\theta)$-framework extension to the current-quark sector, testing whether Paper V's composite successes lift cleanly to the quark level.
\item Running the cascade shell depths to a common scale and computing $\MSbar$-scheme residuals.
\item Structural interpretation of the specific integer-shell decompositions ($104 = ?$, $102 = ?$, etc.) beyond the post-hoc factorisations of Section~\ref{sec:placements}.
\end{enumerate}

\subsection{Recovery Position and Conclusion}
This paper carries the cascade derivation of the Standard Model's six individual quark masses, completing Paper V's thirteen-mass correspondence. The six quarks sit at cascade shell depths $N(m_q) \in [80, 105]$, each associated with an integer $N_*$ that carries a cascade-structural interpretation (\texttt{cascade-masses.md §E3.1}). Three quarks ($u, s, c$) sit within $|\Delta| \le 0.26$ of their integer; three ($d, b, t$) sit at larger offsets. The cascade-derived shell integers include:
\begin{itemize}
\item $s$ and $\mu$ both at shell $96 = 24 \cdot 4$ ($D_4$-root $\times$ $H_4$-eigenvalue-multiplicity); muon precision $0.01\%$.
\item $c, p, n$ all at shell $91$; charm precision $\sim 7\%$.
\item $t$ and $H$ both at shell $81 = 3^4$; top precision $\sim 13\%$.
\end{itemize}
VFD does not replace the Standard Model quark sector. It derives the integer-shell structure from cascade foundations F1 and F3~\citep{SmartXXXVI}, placing each quark on the same $\varphi$-shell scale as the leptons and hadrons, with structural-integer readings documented in~\texttt{cascade-masses.md}. The $\MSbar$ convention-choice remains an external input (one per quark: the PDG-chosen scale); everything else is cascade-derived.

The principal open item is the confinement-to-parton matching layer, which translates the cascade's composite-level $E(\theta)$-framework (Paper V) into current-quark spectral content. That layer is the subject of Paper XXXIX (electroweak dynamics). A second open item is whether the larger offsets ($d, b, t$) reflect real cascade-structural corrections or indicate scheme / running effects not yet modelled at the cascade level.

%==================================================================
\bibliographystyle{plain}
\begin{thebibliography}{99}

\bibitem{RecoveryManifest} L. Smart, \emph{VFD Recovery Manifest}, VFD Programme, 2026. \texttt{docs/recovery-manifest.md}.

\bibitem{CascadeMasses} L. Smart, \emph{Cascade Mass Spectrum: Integer $\varphi$-Shells}, cascade-derivation source, VFD Programme, 2024--2026. \texttt{papers/cascade-derivation/cascade-masses.md}.

\bibitem{SmartI} L. Smart, \emph{The 600-Cell Ansatz}, Paper I, VFD Programme, 2025--2026.
\bibitem{SmartV} L. Smart, \emph{Particle Mass Derivation from the 600-Cell}, Paper V, VFD Programme, 2026.
\bibitem{SmartXXXII} L. Smart, \emph{Hadron Charge Radii as Spectral Coherence Lengths from 600-Cell Closure Geometry}, Paper XXXII, VFD Programme, 2026.
\bibitem{SmartXXXVI} L. Smart, \emph{Cascade Foundations: The Eight Foundations F1--F8 Underlying the VFD Programme}, Paper XXXVI, VFD Programme, 2026.
\bibitem{SmartXXXVII} L. Smart, \emph{Neutrino Sector in $H_4$}, Paper XXXVII, VFD Programme, 2026.
\bibitem{PDG2024} R. L. Workman et al.\ (Particle Data Group), \emph{Review of Particle Physics}, Prog.\ Theor.\ Exp.\ Phys.\ 2024, 083C01 (2024).

\end{thebibliography}

\end{document}
